% Generated draft: requires mathematical and source review.
\documentclass[11pt]{article}
\usepackage[margin=25mm]{geometry}
\usepackage{fontspec}
\setmainfont{Latin Modern Roman}
\usepackage{amsmath,amssymb,mathtools,unicode-math}
\setmathfont{Latin Modern Math}
\usepackage{enumitem,needspace,xurl,hyperref}
\setlength{\parindent}{0pt}
\setlength{\parskip}{5pt}
\setlength{\emergencystretch}{3em}
\newcommand{\N}{\mathbb{N}}
\newcommand{\Z}{\mathbb{Z}}
\newcommand{\Q}{\mathbb{Q}}
\newcommand{\R}{\mathbb{R}}
\newcommand{\C}{\mathbb{C}}
\newcommand{\F}{\mathbb{F}}
\newcommand{\A}{\mathbb{A}}
\newcommand{\PP}{\mathbb{P}}
\newcommand{\E}{\mathbb{E}}
\newcommand{\eps}{\varepsilon}
\newcommand{\dd}{\,\mathrm{d}}
\begin{document}
\section*{UnsolvedMath ID: 20000960}
\textbf{Source problem code:} catalogue entry
\subsection{Definitions and mathematical statement}
Let a two-coloring be any map $\chi:\R^2\to\{0,1\}$, without an unstated measurability assumption. A finite set occurs monochromatically if an isometric, hence congruent, copy has constant color; rescaling is not allowed unless explicitly specified. The source conjectures that every three-point set occurs, except possibly one equilateral congruence class specific to the coloring. Thus the exceptional side length may depend on $\chi$, and there is no proposed universal exceptional scale.

For each coloring define its exceptional equilateral spectrum
\[
 E(\chi)=\{s>0:\text{no monochromatic equilateral triangle of side }s\text{ exists}\}.
\]
The aid reports, using a theorem abbreviated EGMRSS in the summary, that the nonequilateral-triple conjecture is equivalent to $|E(\chi)|\leq1$ for every two-coloring. The exact side-length theorem and its hypotheses are not supplied, so that equivalence is retained as an attributed reduction, not reconstructed without its cited input.

For a precise plane-stripe special case, fix $h>0$ and a unit vector $u\in\R^2$, and set
\[
 \chi_{h,u}(x)=\left\lfloor\frac{\langle x,u\rangle}{h}\right\rfloor\pmod2.
\]
The stripes are half open. The aid's exact spectrum theorem states
\[
 E(\chi_{h,u})=\{2h/\sqrt3\}.
\]
Equivalently, an equilateral triangle is avoided exactly when its altitude equals $h$; every other positive side length has a monochromatic congruent copy. In particular $h=\sqrt3/2$ gives the source's unit-equilateral counterexample. This plane assertion is not transferred to higher-dimensional slab colorings, where additional orientations change the geometry.

The reported boundary-sensitive criterion uses projected vertex positions modulo $2h$. For a fixed orientation of a triangle, take the distinct residues of its three heights on the circle of circumference $2h$ and let $g_{\max}$ be their largest cyclic gap. A translation making that oriented triple monochromatic exists exactly when
\[
 g_{\max}>h.
\]
The strict sign matters: an arc of length exactly $h$ with both endpoints required does not fit inside a half-open stripe. Rotations are then varied in determining the equilateral spectrum.

The original note separately reports that for a fixed set $T$ of three collinear points there is a sixteen-coloring of $\mathbb E^N$ avoiding every congruent copy of $T$. Its dimension quantifier and full reference are not supplied in the literal record; it is a many-color avoidance statement, not a two-color counterexample to the plane conjecture. The stripe-spectrum result does not solve the arbitrary-coloring conjecture.
\subsection{Short English statement}
The exception concerns a fixed equilateral side length for each coloring, not an arbitrary rescaling. The supplied theorem finds exactly that one excluded length for alternating half-open stripes. A reported general equivalence uses an incompletely cited side-length theorem, while the full two-color conjecture remains unresolved here.
\subsection{Sources}
\begin{itemize}
\item Local UnsolvedMath registry, permanent ID $20000960$, source code \texttt{catalogue entry}, supplied full source record and research summary.
\item AIM workshop \emph{Graph Ramsey theory}, Euclidean Ramsey sets, Conjecture 38, supplied Overleaf source \url{https://www.overleaf.com/read/mnvcscjjysvg}; canonical local record \texttt{aim-combinatorics-notes.json}, \texttt{notes[84]}.
\item The summary invokes an ``EGMRSS side-length theorem'' without a full bibliographic entry or statement. The original sixteen-color collinear-set note likewise lacks a citation; the original literature field is empty.
\item Research attempt 1: partial result; problem open at the recorded run, registry status partially solved; supplied rigor label ``proved'', novelty confidence ``low''. The machine-generated aid is not peer reviewed and requires independent verification before citation.
\end{itemize}
\end{document}
